\chapter{System Overview}
\label{chap:system_overview}

This chapter details the architectural design and modelling of Vital Guardian.
The system is built upon a modular, multi-layered, edge-computing architecture
chosen to prioritise patient privacy and real-time performance.

% =============================================================
\section{Architectural Design}
\label{sec:architecture}

The architecture decomposes the system into four distinct subsystems that
collaborate to transform raw sensor data into actionable clinical alerts.

\begin{itemize}
    \item \textbf{Input Layer:} Captures raw video and audio streams from a
    standard camera and microphone connected to the edge device.

    \item \textbf{Processing Layer:} Runs the Visual Guardian (YOLOv11n,
    MoViNet-A2, and MediaPipe Pose) and the Auditory Watchdog (Silero VAD,
    YAMNet, and Faster-Whisper) in parallel on the edge device.

    \item \textbf{Reasoning Layer:} The Cognitive Core uses Gemini language
    models to fuse structured event payloads from the Processing Layer, determine
    contextual severity, and generate natural language alert narratives.

    \item \textbf{Output Layer:} A FastAPI and WebSocket-powered dashboard
    displays tiered alerts and a live video feed to authorised nursing staff.
\end{itemize}

\begin{figure}[H]
    \centering
    \includegraphics[width=0.9\textwidth]{images/arch_diagram.png}
    \caption{High-Level System Architecture of Vital Guardian. Hardware sensors
    (microphone and camera) operate externally to the software architecture boundary.}
    \label{fig:architecture}
\end{figure}

% =============================================================
\section{Design Models}
\label{sec:design_models}

The following object-oriented design models are presented to detail the static
structure and dynamic behaviour of the system.

\subsection{Activity Diagram}

The activity diagram illustrates the complete patient monitoring workflow, from
raw sensor capture through parallel event detection to the final generation of
a multi-modal alert on the nursing dashboard.

\begin{figure}[H]
    \centering
    \includegraphics[width=0.6\textwidth]{images/activity_diagram.png}
    \caption{Activity Diagram for the Complete Patient Monitoring Workflow}
    \label{fig:activity}
\end{figure}

\subsection{Class Diagram}

The class diagram outlines the static software structure. A
\texttt{MonitoringSession} controller manages the \texttt{VisualGuardian} and
\texttt{AuditoryWatchdog} instances and forwards their event outputs to the
\texttt{CognitiveCore}. The diagram implements an inheritance hierarchy in which
\texttt{Patient} and \texttt{NursingStaff} generalise to a common \texttt{Person}
base class.

\begin{figure}[H]
    \centering
    \includegraphics[width=0.9\textwidth]{images/class_diagram.png}
    \caption{Class Diagram Illustrating Main Software Components and Hierarchies}
    \label{fig:class_diagram}
\end{figure}

\subsection{Sequence Diagram}

The sequence diagram details the parallel execution of event detection during a
monitoring cycle. The \texttt{MonitoringSession} triggers
\texttt{createVisionEvent} and \texttt{createAudioEvent} concurrently, collects
the resulting payloads, and forwards them to the \texttt{CognitiveCore} for
verification and enrichment.

\begin{figure}[H]
    \centering
    \includegraphics[width=1.0\textwidth]{images/sequence_diagram.png}
    \caption{Sequence Diagram: Main Monitoring and Alert Generation Cycle}
    \label{fig:sequence}
\end{figure}

\subsection{State Transition Diagram}

This diagram models the lifecycle of a \texttt{MonitoringSession}, transitioning
between states including \texttt{Gathering Events}, \texttt{Evaluating},
\texttt{Alert Triggered}, and \texttt{Alert Acknowledged}.

\begin{figure}[H]
    \centering
    \includegraphics[width=1.0\textwidth]{images/state_diagram.png}
    \caption{State Transition Diagram for the MonitoringSession Object}
    \label{fig:state_diagram}
\end{figure}

% =============================================================
\section{Data Design}
\label{sec:data_design}

The system data design adheres to a local-only processing principle. The two
primary structured data objects are as follows.

\begin{itemize}
    \item \textbf{The Event Object:} An atomic unit of information generated by
    a sensory module. It contains attributes including \texttt{eventID},
    \texttt{eventType} (e.g., \texttt{visual\_fall}), \texttt{timestamp},
    and \texttt{confidence}.

    \item \textbf{The Alert Object:} The verified output produced by the
    Cognitive Core. It contains \texttt{alertID}, \texttt{priority}
    (CRITICAL, HIGH, MEDIUM, or LOW), \texttt{alertMessage} (a human-readable
    narrative), and a list of \texttt{contributingEvents} for full traceability.
\end{itemize}

All data is processed in-memory for real-time operations. Verified alerts are
persisted to a local PostgreSQL database via SQLAlchemy for auditing and review
by nursing staff. No raw video or audio data is ever written to disk or
transmitted externally.

Figure \ref{fig:dfd} illustrates the flow of these data objects from the raw
sensor inputs through the processing modules to the nursing dashboard.

\begin{figure}[H]
    \centering
    \includegraphics[width=1.0\textwidth]{images/dfd_diagram.png}
    \caption{Data Flow Diagram for Vital Guardian}
    \label{fig:dfd}
\end{figure}

% =============================================================
\section{Domain Model}
\label{sec:domain_model}

The domain model illustrates the conceptual entities within the system and their
relationships. Key entities include the \texttt{Patient}, \texttt{NursingStaff},
\texttt{MonitoringSession}, and the three core software components:
\texttt{VisualGuardian}, \texttt{AuditoryWatchdog}, and \texttt{CognitiveCore}.

\begin{figure}[H]
    \centering
    \includegraphics[width=0.8\textwidth]{images/domain_model.png}
    \caption{Conceptual Domain Model for Vital Guardian}
    \label{fig:domain_model}
\end{figure}
